\chapter{Остров в цифрах}

\factline{Площадь 1560 км². Население ок.~855~000. Высшая точка --- Пико-де-лас-Ньевес, 1949~м.}

\place{Гран-Канария} --- третий по~величине из~Канарских островов,
почти круглой формы около 50~км в~поперечнике. Его называют
\textit{«континентом в~миниатюре»}: на~коротком переезде
с~севера на~юг путешественник проходит четыре климатические
зоны --- от~влажных лавровых склонов к~сухим полупустынным дюнам.

\marginfact{Координаты\\центра острова:\\27°57$'$\,N,\\15°36$'$\,W}

\section{Геология}

Остров имеет чисто вулканическое происхождение; возраст старейших
базальтовых пород около 14\,млн лет. В~центре возвышается древний
кальдерный массив, окружённый радиальной системой глубоких
\textit{барранко} (ущелий). Характерные образования ---
\sight{Роке-Нубло} (вулканическая \textit{игла} высотой 80~м)
и~\sightt{Дюны Маспаломаса}, единственное в~Европе образование
подвижных песков такого масштаба.

\section{Климат}

Благодаря пассатам и~холодному Канарскому течению климат мягкий
круглый год. Север --- зелёный, с~частыми низкими облаками
(\textit{panza de burro}, «ослиное брюхо»); юг --- сухой и~солнечный.

\begin{center}\small
\begin{tabular}{lcccc}
\toprule
Месяц      & Янв & Апр & Июл & Окт \\
\midrule
Воздух, ° C (Лас-Пальмас) & 18 & 19 & 24 & 23 \\
Вода, ° C                 & 19 & 19 & 22 & 23 \\
Осадки, мм                      & 30 & 8  & 0  & 15 \\
\bottomrule
\end{tabular}
\end{center}

\section{Растительность}

На~северных склонах сохраняются остатки \textit{лаврисильвы} ---
субтропического туманного леса. Выше 800~м их~сменяет
\place{канарская сосна} (\textit{Pinus canariensis}),
удивительно живучая порода, выдерживающая лесные пожары.
Прибрежная полоса занята суккулентами, молочаями
(\textit{cardón}) и~драконовыми деревьями (\textit{drago}).

\begin{detour}
Любителям ботаники --- \sight{Jardín Canario} в~Таферо, крупнейший
ботанический сад Испании, посвящённый эндемикам
Макаронезии; вход свободный.
\end{detour}
